\chapter{Industrialisation : du diagnostic à la solution déployable}
\label{chap:indus}

Ce chapitre regroupe les développements qui font passer PrognoSense d'un démonstrateur
d'analyse à une solution \emph{déployable} et \emph{digne de confiance} sur le terrain.
C'est ici que se concentre la valeur différenciante du projet.

\section{Sévérité normalisée ISO 10816/20816}
La plateforme exprime la sévérité dans le langage des ingénieurs de fiabilité : vitesse
vibratoire efficace en mm/s et zones A/B/C/D par classe de machine (cf.
chapitre~\ref{chap:pipeline}). La zone ISO est calculée à l'ingestion, injectée dans le
diagnostic et affichée dans les ordres de travail. Cela rend chaque diagnostic
\textbf{comparable et défendable}, à la différence d'un indice propriétaire.

\section{Connectivité terrain : du capteur à l'ingestion}
\subsection{Couche d'ingestion universelle}
Le point d'entrée \texttt{POST /api/ingest/signal} accepte une forme d'onde
d'accélération réelle, quelle qu'en soit la source, et la fait traverser tout le pipeline :
extraction d'indicateurs $\rightarrow$ sévérité ISO $\rightarrow$ anomalie vs baseline
machine $\rightarrow$ diagnostic spectral $\rightarrow$ indice de santé $\rightarrow$
persistance et mise à jour de la flotte. Cette abstraction \emph{vendor-agnostic} est ce qui
transforme la plateforme en système connectable au parc existant.

\subsection{OPC-UA}
Un connecteur OPC-UA (bibliothèque \texttt{asyncua}) lit périodiquement un n\oe ud de mesure
sur un serveur OPC-UA --- le protocole standard de l'automatisme industriel. La chaîne
complète a été validée localement à l'aide d'un serveur OPC-UA de démonstration : la machine
virtuelle \texttt{POMPE\_OPCUA} a été correctement lue puis ingérée (états persistés,
baseline créé), prouvant le fonctionnement de bout en bout \emph{capteur} $\rightarrow$
\emph{ingestion}.

\subsection{MQTT}
Un connecteur MQTT (\texttt{paho-mqtt}) souscrit à un topic IIoT et transfère chaque message
vers la même ingestion. Les deux connecteurs sont \emph{activés sur configuration}
(variables d'environnement), sans impact lorsqu'ils ne sont pas utilisés.

\section{Apprentissage du comportement sain par machine}
Comme détaillé au chapitre~\ref{chap:anomalie}, le module \texttt{MachineBaseline} apprend
la signature saine \emph{propre à chaque machine}, sans nécessiter d'historique de pannes.
L'anomalie est mesurée relativement à cette référence : validation à $\approx 5\,\%$ sur un
signal sain et $\approx 96\,\%$ sur un signal dégradé, avec désignation des indicateurs
contributeurs.

\section{Boucle fermée : de l'alerte à l'ordre de travail}
PrognoSense ferme la boucle de maintenance. Lorsqu'une machine atteint un état critique
(zone ISO D ou indice de santé bas), un \textbf{ordre de travail} est créé automatiquement,
avec priorité (P1/P2), machine, défaut et action recommandée, en évitant les doublons (un
seul ordre ouvert par machine). Cet ordre peut être \textbf{poussé vers la GMAO} de
l'entreprise (SAP PM, IBM Maximo, Infor) via un connecteur webhook configurable. L'alerte
devient ainsi une intervention planifiée et tracée. Des notifications par e-mail (SMTP)
complètent le dispositif.

\section{Indicateurs de fiabilité : mesurés, pas estimés}
\subsection{Un choix méthodologique assumé}
Un point central distingue PrognoSense de nombreux prototypes : les indicateurs de fiabilité
ne sont pas \emph{estimés} à partir de la simulation, mais \textbf{mesurés sur l'historique
réel des interventions}, conformément à leur définition normative :
\begin{align}
\text{MTBF} &= \frac{\text{temps de fonctionnement cumulé}}{\text{nombre de pannes réelles}} &
\text{MTTR} &= \frac{\text{temps total de réparation}}{\text{nombre de réparations}} \\
\text{Disponibilité} &= \frac{\text{MTBF}}{\text{MTBF}+\text{MTTR}}\times 100 & & \notag
\end{align}
Un MTBF ne se déduit pas d'un capteur : il se \emph{constate} sur les pannes effectives.
C'est pourquoi la plateforme propose une saisie des événements de maintenance (le \og carnet
de bord \fg) qui alimente ces calculs. Tant qu'aucune donnée n'est saisie, la plateforme
\emph{n'affiche rien} plutôt qu'un chiffre fabriqué --- gage d'honnêteté scientifique.

\subsection{Retour sur investissement}
À partir des interventions, la plateforme estime les arrêts évités et les économies
réalisées (paramètres ajustables : coût horaire de production, durée d'arrêt évitée),
fournissant un argumentaire de ROI exploitable par la direction.

\subsection{Efficacité prédictive et taux de fausses alarmes}
En croisant les alertes émises avec les pannes réelles (fenêtre d'anticipation de 48\,h), la
plateforme calcule la \textbf{précision}, le \textbf{rappel} et le \textbf{lead-time}
(délai d'anticipation moyen). Surtout, un mécanisme de \textbf{validation par l'expert}
(verdict vrai/faux positif sur chaque alerte) permet de mesurer le \emph{vrai} taux de
fausses alarmes --- l'indicateur de confiance le plus exigé par les acheteurs de solutions
PdM. Cette boucle \emph{human-in-the-loop} est, à la connaissance de l'auteur, rarement
présente dans les travaux comparables.

\section{Copilot ancré sur les connaissances (RAG)}
Un assistant conversationnel, propulsé par le modèle de langage \textbf{Mistral}, répond en
langage naturel et dispose du contexte temps réel de la flotte. Il est \textbf{ancré} sur une
base de connaissances (normes ISO, guide des signatures de défaut, interprétation des
indicateurs) via une recherche de similarité \textbf{TF-IDF} (RAG). Ses réponses
\emph{citent les sources documentaires}, ce qui les rend précises et \textbf{traçables}.
Exemple validé : à la question \og quel défaut provoque un pic à 2$\times$ la fréquence de
rotation ?\fg, le copilot répond \og désalignement \fg{} en s'appuyant sur le guide des
défauts, sources à l'appui. Le copilot démocratise ainsi l'expertise d'un analyste vibratoire
auprès des techniciens.

\section{Confiance et gouvernance des modèles}
\subsection{Détection de dérive (concept drift)}
À chaque prédiction, un \texttt{DriftDetector} compare la distribution des données de
production à la distribution d'entraînement par test de \textbf{Kolmogorov-Smirnov}. Une
dérive est signalée lorsque plus de 30\,\% des features testées présentent une p-value
$< 0{,}05$, avec une sévérité critique si $\min(p)<0{,}001$. La plateforme recommande alors
un réentraînement.

\subsection{Calibration des probabilités}
La calibration (isotonique, \texttt{CalibratedClassifierCV}) corrige le biais de confiance :
sans elle, un modèle annonçant \og 94\,\% \fg{} peut n'être correct que dans 75\,\% des cas.
Un diagramme de fiabilité et l'erreur de calibration espérée (ECE) sont calculés ; un modèle
est jugé bien calibré si ECE $< 0{,}05$.

\subsection{Versioning et retour arrière}
Chaque réentraînement crée une \textbf{version} archivée du modèle (métriques, date,
déclencheur). En cas de régression, un \textbf{rollback} en un clic réactive une version
antérieure. Cette discipline MLOps garantit qu'aucune mise à jour ne dégrade durablement le
service.

\subsection{Journal d'audit}
Chaque décision et chaque alerte est consignée et horodatée (prédiction, machine, confiance),
avec persistance sur disque survivant aux redémarrages --- exigence de conformité (esprit
ISO~13373) et de traçabilité.

\section{Passage à l'échelle (time-series)}
Pour un parc réel, la plateforme est \textbf{prête à basculer sur TimescaleDB} (PostgreSQL +
extension time-series) sans modification du code : seule la variable de connexion change.
Une lecture \textbf{agrégée par downsampling} (validée : plusieurs centaines de points bruts
condensés en tranches horaires) et une \textbf{politique de rétention} sont fournies, ainsi
qu'un \texttt{docker-compose} et une documentation de mise à l'échelle.
